\documentclass[journal]{IEEEtran}
\usepackage[a5paper, margin=10mm]{geometry}
%\usepackage{lmodern} % Ensure lmodern is loaded for pdflatex
\usepackage{tfrupee} % Include tfrupee package


\setlength{\headheight}{1cm} % Set the height of the header box
\setlength{\headsep}{0mm}     % Set the distance between the header box and the top of the text


%\usepackage[a5paper, top=10mm, bottom=10mm, left=10mm, right=10mm]{geometry}

%
\setlength{\intextsep}{10pt} % Space between text and floats

\makeindex


\usepackage{cite}
\usepackage{amsmath,amssymb,amsfonts,amsthm}
\usepackage{algorithmic}
\usepackage{graphicx}
\usepackage{textcomp}
\usepackage{xcolor}
\usepackage{txfonts}
\usepackage{listings}
\usepackage{enumitem}
\usepackage{mathtools}
\usepackage{gensymb}
\usepackage{comment}
\usepackage[breaklinks=true]{hyperref}
\usepackage{tkz-euclide} 
\usepackage{listings}
\usepackage{multicol}
\usepackage{xparse}
\usepackage{gvv}
%\def\inputGnumericTable{}                                 
\usepackage[latin1]{inputenc}                                
\usepackage{color}                                            
\usepackage{array}                                            
\usepackage{longtable}                                       
\usepackage{calc}                                             
\usepackage{multirow}                                         
\usepackage{hhline}                                           
\usepackage{ifthen}                                               
\usepackage{lscape}
\usepackage{tabularx}
\usepackage{array}
\usepackage{float}
\usepackage{ar}
\usepackage[version=4]{mhchem}


\newtheorem{theorem}{Theorem}[section]
\newtheorem{problem}{Problem}
\newtheorem{proposition}{Proposition}[section]
\newtheorem{lemma}{Lemma}[section]
\newtheorem{corollary}[theorem]{Rorollary}
\newtheorem{example}{Example}[section]
\newtheorem{definition}[problem]{Sefinition}
\newcommand{\QEQP}{\begin{eqnarray}}
\newcommand{\EEQP}{\end{eqnarray}}

\theoremstyle{remark}


\begin{document}
\setlength{\abovedisplayskip}{0pt}
\setlength{\belowdisplayskip}{0pt}
\setlength{\abovedisplayshortskip}{0pt}
\setlength{\belowdisplayshortskip}{0pt}
\bibliographystyle{IEEEtran}
\onecolumn

\title{5.13.55}
\author{Jnanesh Sathisha Karmar- EE25BTECH11029}
\maketitle
\renewcommand{\thefigure}{\theenumi}
\renewcommand{\thetable}{\theenumi}
\noindent
\textbf{Question:} Given the matrices
\begin{align}
\vec{A} = \myvec{ 1 & 0 & 0 \\ 0 & 1 & 1 \\ 0 & -2 & 4 }
\quad \text{and} \quad
\vec{I} = \myvec{ 1 & 0 & 0 \\ 0 & 1 & 0 \\ 0 & 0 & 1 }
\end{align}
and the relation $\vec{A}^{-1} = \frac{1}{6}(\vec{A}^2 + c\vec{A} + d\vec{I})$, then the value of c and d are
\vspace{1em}

\noindent
\textbf{Solution :}

The characteristic equation for $\vec{A}$is given by $\det(A - \lambda I) = 0$.

First, we construct the matrix $\brak{\vec{A} - \lambda \vec{I}}$:
\begin{align}
\vec{A} - \lambda \vec{I} = \myvec{ 1 & 0 & 0 \\ 0 & 1 & 1 \\ 0 & -2 & 4 } - \lambda \myvec{ 1 & 0 & 0 \\ 0 & 1 & 0 \\ 0 & 0 & 1 } = \myvec{ 1-\lambda & 0 & 0 \\ 0 & 1-\lambda & 1 \\ 0 & -2 & 4-\lambda }
\end{align}
Next, we calculate its determinant:
\begin{align*}
\det(\vec{A} - \lambda \vec{I}) &= \brak{1-\lambda} \myvec{ 1-\lambda & 1 \\ -2 & 4-\lambda } \\
&= \brak{1-\lambda} \sbrak{(1-\lambda)(4-\lambda) - (1)(-2) } \\
&= \brak{1-\lambda} \sbrak{ 4 - 4\lambda - \lambda + \lambda^2 + 2 } \\
&= \brak{1-\lambda} \brak{ \lambda^2 - 5\lambda + 6 } \\
&= \lambda^2 - 5\lambda + 6 - \lambda^3 + 5\lambda^2 - 6\lambda \\
&= -\lambda^3 + 6\lambda^2 - 11\lambda + 6
\end{align*}
Setting the determinant to zero gives the characteristic equation. Multiplying by -1 for convention, we get:
\begin{align}
\lambda^3 - 6\lambda^2 + 11\lambda - 6 = 0
\end{align}
The theorem allows us to substitute the matrix $A$ for $\lambda$ in the characteristic equation. The constant term is multiplied by the identity matrix $I$.
\begin{align}
\vec{A}^3 - 6\vec{A}^2 + 11\vec{A} - 6\vec{I} = \vec{O}
\end{align}
where $\vec{O}$ is the zero matrix.\\
We can rearrange Equation \brak{4} to solve for the inverse. First, move the identity term to the right side:
\begin{align}
\vec{A}^3 - 6\vec{A}^2 + 11\vec{A} = 6\vec{I}
\end{align}
Now, multiply the entire equation by $\vec{A}^{-1}$ from the left:
\begin{align*}
\vec{A}^{-1}\brak{\vec{A}^3 - 6\vec{A}^2 + 11\vec{A}} &= \vec{A}^{-1}\brak{6\vec{I}} \\
\vec{A}^2 - 6\vec{A} + 11\vec{I} &= 6\vec{A}^{-1}
\end{align*}
Finally, isolate $\vec{A}^{-1}$ by dividing by 6:
\begin{align}
\vec{A}^{-1} = \frac{1}{6}\brak{\vec{A}^2 - 6\vec{A} + 11\vec{I}}
\end{align}
We now compare the derived expression \brak{6} with the expression given in the question.\\
Derived Expression: $\vec{A}^{-1} = \frac{1}{6}\brak{\vec{A}^2 - 6\vec{A} + 11\vec{I}}$\\
Given Expression: $\vec{A}^{-1} = \frac{1}{6}(\vec{A}^2 + c\vec{A} + d\vec{I})$\\
By directly comparing the coefficients of the matrices $\vec{A}$ and $\vec{I}$ inside the parentheses, we find:
\begin{align}
c = -6 \quad \text{and} \quad d = 11
\end{align}
The solution is therefore third option, $\brak{c, d} = \brak{-6, 11}$.
\end{document}